% Auto-generated by scripts/paper_parsing.py from arXiv-2405.15459v2.
% Sources parsed: assembler.tex, math_commands.tex.
% Intra-paper duplicate names (86): 0, 1, Eb, N, R, aa, bb, cA, cB, cC, cD, cE, cF, cG, cH, cJ, cK, cL, cM, cN, cO, cP, cQ, cR, cS, cT, cU, cV, cW, cX, cY, cZ, cc, dd, ee, ff, hh, ii, jj, lll, mA, mB, mC, mD, mE, mF, mG, mH, mI, mJ, mK, mL, mLambda, mM, mN, mO, mP, mQ, mR, mS, mT, mU, mV, mW, mX, mY, mZ, mm, nn, oo, pp, qq, refEQ, refGE, refID, refLE, rr, ss, tt, uu, vec, vv, ww, xx, yy, zz

% --- Math operators ---
% from math_commands.tex:15
\DeclareMathOperator*{\Circ}{\bigcirc}
% from math_commands.tex:280
\DeclareMathOperator{\Cov}{Cov}
% from math_commands.tex:281
\DeclareMathOperator{\Span}{Span}
% from math_commands.tex:282
\DeclareMathOperator{\Relu}{ReLU}

% --- Commands ---
% from assembler.tex:29
\renewcommand{\vec}[1]{\bm{#1}}
% from assembler.tex:30
\newcommand{\change}[1]{\textcolor{purple}{ #1}}
% from assembler.tex:63
\newcommand{\otherlabel}[2]{\protected@edef\@currentlabel{#2}\label{#1}}
% from math_commands.tex:5
\newcommand
% from math_commands.tex:7
\newcommand{\dP}{\mathbb{P}}
% from math_commands.tex:16
\newcommand{\Ea}[1]{\E\left[#1\right]}
% from math_commands.tex:17
\newcommand{\Eb}[2]{\E_{#1}\left[#2\right]}
% from math_commands.tex:18
\newcommand{\bG}{{\boldsymbol{G}}}
% from math_commands.tex:19
\newcommand{\bT}{{\boldsymbol{T}}}
% from math_commands.tex:20
\newcommand{\bOne}{{\boldsymbol{1}}}
% from math_commands.tex:21
\newcommand{\bR}{{\boldsymbol{R}}}
% from math_commands.tex:22
\newcommand{\bhR}{{\hat{\boldsymbol{R}}}}
% from math_commands.tex:23
\newcommand{\bg}{{\boldsymbol{g}}}
% from math_commands.tex:24
\newcommand{\bH}{{\boldsymbol{H}}}
% from math_commands.tex:25
\newcommand{\bA}{{\boldsymbol{A}}}
% from math_commands.tex:26
\newcommand{\bC}{{\boldsymbol{C}}}
% from math_commands.tex:27
\newcommand{\bS}{{\boldsymbol{S}}}
% from math_commands.tex:28
\newcommand{\bhS}{{\hat{\boldsymbol{S}}}}
% from math_commands.tex:29
\newcommand{\bU}{{\boldsymbol{U}}}
% from math_commands.tex:30
\newcommand{\bhU}{{\hat{\boldsymbol{U}}}}
% from math_commands.tex:31
\newcommand{\bh}{{\boldsymbol{h}}}
% from math_commands.tex:32
\newcommand{\bW}{{{\boldsymbol{{W}}}}}
% from math_commands.tex:33
\newcommand{\bw}{{\boldsymbol{w}}}
% from math_commands.tex:34
\newcommand{\bQ}{{\boldsymbol{Q}}}
% from math_commands.tex:35
\newcommand{\bV}{{\boldsymbol{V}}}
% from math_commands.tex:36
\newcommand{\bsQ}{{\boldsymbol{\mathsf Q}}}
% from math_commands.tex:37
\newcommand{\bsS}{{\boldsymbol{\mathsf S}}}
% from math_commands.tex:38
\newcommand{\bM}{{\boldsymbol{m}}}
% from math_commands.tex:39
\newcommand{\bhQ}{{\hat{\boldsymbol{Q}}}}
% from math_commands.tex:40
\newcommand{\bhV}{{\hat{\boldsymbol{V}}}}
% from math_commands.tex:41
\newcommand{\bhM}{{\hat{\boldsymbol{m}}}}
% from math_commands.tex:42
\newcommand{\bMM}{{{\boldsymbol{M}}}}
% from math_commands.tex:43
\newcommand{\bhMM}{{\hat{\boldsymbol{M}}}}
% from math_commands.tex:44
\newcommand{\bbb}{{\boldsymbol{b}}}
% from math_commands.tex:45
\newcommand{\bv}{{\boldsymbol{v}}}
% from math_commands.tex:46
\newcommand{\ba}{{\boldsymbol{a}}}
% from math_commands.tex:47
\newcommand{\bbB}{{\boldsymbol{B}}}
% from math_commands.tex:48
\newcommand{\bF}{{\boldsymbol{F}}}
% from math_commands.tex:49
\newcommand{\bZ}{{\boldsymbol{Z}}}
% from math_commands.tex:50
\newcommand{\bY}{{\boldsymbol{Y}}}
% from math_commands.tex:51
\newcommand{\btheta}{{\boldsymbol{\theta}}}
% from math_commands.tex:52
\newcommand{\bTheta}{{\boldsymbol{\Theta}}}
% from math_commands.tex:53
\newcommand{\bSigma}{{\boldsymbol{\Sigma}}}
% from math_commands.tex:54
\newcommand{\bPi}{{\boldsymbol{\Pi}}}
% from math_commands.tex:55
\newcommand{\bDelta}{{\boldsymbol{\Delta}}}
% from math_commands.tex:56
\newcommand{\brho}{{\boldsymbol{\rho}}}
% from math_commands.tex:57
\newcommand{\bphi}{{\boldsymbol{\phi}}}
% from math_commands.tex:58
\newcommand{\bsigma}{{\boldsymbol{\sigma}}}
% from math_commands.tex:59
\newcommand{\blambda}{{\boldsymbol{\lambda}}}
% from math_commands.tex:60
\newcommand{\bomega}{{\boldsymbol{\omega}}}
% from math_commands.tex:61
\newcommand{\bzeta}{{\boldsymbol{\zeta}}}
% from math_commands.tex:62
\newcommand{\bXi}{{\boldsymbol{\Xi}}}
% from math_commands.tex:63
\newcommand{\bmu}{{\boldsymbol{\mu}}}
% from math_commands.tex:64
\newcommand{\bvarphi}{{\boldsymbol{\varphi}}}
% from math_commands.tex:65
\newcommand{\bz}{{\boldsymbol{z}}}
% from math_commands.tex:66
\newcommand{\bx}{{\boldsymbol{x}}}
% from math_commands.tex:67
\newcommand{\bn}{{\boldsymbol{n}}}
% from math_commands.tex:68
\newcommand{\bX}{{\boldsymbol{X}}}
% from math_commands.tex:69
\newcommand{\bI}{{\boldsymbol{I}}}
% from math_commands.tex:71
\newcommand{\bB}{{\boldsymbol{B}}}
% from math_commands.tex:72
\newcommand{\by}{{\boldsymbol{y}}}
% from math_commands.tex:73
\newcommand{\bff}{{\boldsymbol{f}}}
% from math_commands.tex:74
\newcommand{\be}{{\boldsymbol{e}}}
% from math_commands.tex:75
\newcommand{\beeta}{{\boldsymbol{\eta}}}
% from math_commands.tex:76
\newcommand{\bxi}{{\boldsymbol{\xi}}}
% from math_commands.tex:77
\newcommand{\bpi}{{\boldsymbol{\pi}}}
% from math_commands.tex:78
\newcommand{\bu}{{\boldsymbol{u}}}
% from math_commands.tex:80
\newcommand{\cS}{\mathcal{S}}
% from math_commands.tex:81
\newcommand{\bgamma}{{\boldsymbol{\gamma}}}
% from math_commands.tex:82
\newcommand{\bGamma}{{\boldsymbol{\Gamma}}}
% from math_commands.tex:83
\newcommand{\balpha}{{\boldsymbol{\alpha}}}
% from math_commands.tex:84
\newcommand{\bnu}{{\boldsymbol{\nu}}}
% [intra-paper duplicate] \vec — also defined at assembler.tex:29
% from math_commands.tex:85
\renewcommand{\vec}[1]{\boldsymbol{#1}}
% from math_commands.tex:86
\providecommand{\refLE}[1]{\ensuremath{\stackrel{(\ref{#1})}{\leq}}}
% from math_commands.tex:87
\providecommand{\refEQ}[1]{\ensuremath{\stackrel{(\ref{#1})}{=}}}
% from math_commands.tex:88
\providecommand{\refGE}[1]{\ensuremath{\stackrel{(\ref{#1})}{\geq}}}
% from math_commands.tex:89
\providecommand{\refID}[1]{\ensuremath{\stackrel{(\ref{#1})}{\equiv}}}
% from math_commands.tex:91
\providecommand{\mgeq}{\succeq}
% from math_commands.tex:92
\providecommand{\mleq}{\preceq}
% from math_commands.tex:94
\providecommand{\divides}{\mid}
% from math_commands.tex:95
\providecommand{\tsum}{\textstyle\sum}
% from math_commands.tex:97
\providecommand{\R}{\mathbb{R}}
% from math_commands.tex:98
\providecommand{\real}{\mathbb{R}}
% from math_commands.tex:99
\providecommand{\N}{\mathbb{N}}
% from math_commands.tex:103
\def\sign{\@ifnextchar*{\@sgnargscaled}{\@ifnextchar[{\sgnargscaleas}{\@ifnextchar{\bgroup}{\@sgnarg}{\sgn} }}}
% from math_commands.tex:104
\def\@sgnarg#1{\sgn\rbr{#1}}
% from math_commands.tex:105
\def\@sgnargscaled#1{\sgn\rbr*{#1}}
% from math_commands.tex:106
\def\@sgnargscaleas[#1]#2{\sgn\rbr[#1]{#2}}
% from math_commands.tex:110
\providecommand{\0}{\bm{0}}
% from math_commands.tex:111
\providecommand{\1}{\bm{1}}
% from math_commands.tex:112
\providecommand{\alphav}{\bm{\alpha}}
% from math_commands.tex:113
\renewcommand{\aa}{\bm{a}}
% from math_commands.tex:114
\providecommand{\bb}{\bm{b}}
% from math_commands.tex:115
\providecommand{\cc}{\bm{c}}
% from math_commands.tex:116
\providecommand{\dd}{\bm{d}}
% from math_commands.tex:117
\providecommand{\ee}{\bm{e}}
% from math_commands.tex:118
\providecommand{\ff}{\bm{f}}
% from math_commands.tex:119
\let\ggg\gg
% from math_commands.tex:120
\providecommand{\hh}{\bm{h}}
% from math_commands.tex:121
\providecommand{\ii}{\bm{i}}
% from math_commands.tex:122
\providecommand{\jj}{\bm{j}}
% from math_commands.tex:123
\providecommand{\kk}{\bm{k}}
% from math_commands.tex:124
\let\lll\ll
% from math_commands.tex:126
\providecommand{\mm}{\bm{m}}
% from math_commands.tex:127
\providecommand{\nn}{\bm{n}}
% from math_commands.tex:128
\providecommand{\oo}{\bm{o}}
% from math_commands.tex:129
\providecommand{\pp}{\bm{p}}
% from math_commands.tex:130
\providecommand{\qq}{\bm{q}}
% from math_commands.tex:131
\providecommand{\rr}{\bm{r}}
% from math_commands.tex:132
\let\sss\ss
% from math_commands.tex:133
\renewcommand{\ss}{\bm{s}}
% from math_commands.tex:134
\providecommand{\tt}{\bm{t}}
% from math_commands.tex:135
\providecommand{\uu}{\bm{u}}
% from math_commands.tex:136
\providecommand{\vv}{\bm{v}}
% from math_commands.tex:137
\providecommand{\ww}{\bm{w}}
% from math_commands.tex:138
\providecommand{\xx}{\bm{x}}
% from math_commands.tex:139
\providecommand{\yy}{\bm{y}}
% from math_commands.tex:140
\providecommand{\ddot}{\dot\dot}
% from math_commands.tex:141
\providecommand{\dddot}{\dot\dot}
% from math_commands.tex:142
\providecommand{\zz}{\bm{z}}
% from math_commands.tex:143
\providecommand{\thth}{\bm{\theta}}
% from math_commands.tex:144
\providecommand{\E}{\mathbb{E}}
% from math_commands.tex:146
\providecommand{\txx}{\tilde\xx}
% from math_commands.tex:147
\providecommand{\tgg}{\tilde\gg}
% from math_commands.tex:151
\providecommand{\mA}{\bm{A}}
% from math_commands.tex:152
\providecommand{\mB}{\bm{B}}
% from math_commands.tex:153
\providecommand{\mC}{\bm{C}}
% from math_commands.tex:154
\providecommand{\mD}{\bm{D}}
% from math_commands.tex:155
\providecommand{\mE}{\bm{E}}
% from math_commands.tex:156
\providecommand{\mF}{\bm{F}}
% from math_commands.tex:157
\providecommand{\mG}{\bm{G}}
% from math_commands.tex:158
\providecommand{\mH}{\bm{H}}
% from math_commands.tex:159
\providecommand{\mI}{\bm{I}}
% from math_commands.tex:160
\providecommand{\mJ}{\bm{J}}
% from math_commands.tex:161
\providecommand{\mK}{\bm{K}}
% from math_commands.tex:162
\providecommand{\mL}{\bm{L}}
% from math_commands.tex:163
\providecommand{\mM}{\bm{M}}
% from math_commands.tex:164
\providecommand{\mN}{\bm{N}}
% from math_commands.tex:165
\providecommand{\mO}{\bm{O}}
% from math_commands.tex:166
\providecommand{\mP}{\bm{P}}
% from math_commands.tex:167
\providecommand{\mQ}{\bm{Q}}
% from math_commands.tex:168
\providecommand{\mR}{\bm{R}}
% from math_commands.tex:169
\providecommand{\mS}{\bm{S}}
% from math_commands.tex:170
\providecommand{\mT}{\bm{T}}
% from math_commands.tex:171
\providecommand{\mU}{\bm{U}}
% from math_commands.tex:172
\providecommand{\mV}{\bm{V}}
% from math_commands.tex:173
\providecommand{\mW}{\bm{W}}
% from math_commands.tex:174
\providecommand{\mX}{\bm{X}}
% from math_commands.tex:175
\providecommand{\mY}{\bm{Y}}
% from math_commands.tex:176
\providecommand{\mZ}{\bm{Z}}
% from math_commands.tex:177
\providecommand{\mLambda}{\bm{\Lambda}}
% from math_commands.tex:180
\providecommand{\cA}{\mathcal{A}}
% from math_commands.tex:181
\providecommand{\cB}{\mathcal{B}}
% from math_commands.tex:182
\providecommand{\cC}{\mathcal{C}}
% from math_commands.tex:183
\providecommand{\cD}{\mathcal{D}}
% from math_commands.tex:184
\providecommand{\cE}{\mathcal{E}}
% from math_commands.tex:185
\providecommand{\cF}{\mathcal{F}}
% from math_commands.tex:186
\providecommand{\cG}{\mathcal{G}}
% from math_commands.tex:187
\providecommand{\cH}{\mathcal{H}}
% from math_commands.tex:188
\providecommand{\cII}{\mathcal{H}}
% from math_commands.tex:189
\providecommand{\cJ}{\mathcal{J}}
% from math_commands.tex:190
\providecommand{\cK}{\mathcal{K}}
% from math_commands.tex:191
\providecommand{\cL}{\mathcal{L}}
% from math_commands.tex:192
\providecommand{\cM}{\mathcal{M}}
% from math_commands.tex:193
\providecommand{\cN}{\mathcal{N}}
% from math_commands.tex:194
\providecommand{\cO}{\mathcal{O}}
% from math_commands.tex:195
\providecommand{\cP}{\mathcal{P}}
% from math_commands.tex:196
\providecommand{\cQ}{\mathcal{Q}}
% from math_commands.tex:197
\providecommand{\cR}{\mathcal{R}}
% [intra-paper duplicate] \cS — also defined at math_commands.tex:80
% from math_commands.tex:198
\providecommand{\cS}{\mathcal{S}}
% from math_commands.tex:199
\providecommand{\cT}{\mathcal{T}}
% from math_commands.tex:200
\providecommand{\cU}{\mathcal{U}}
% from math_commands.tex:201
\providecommand{\cV}{\mathcal{V}}
% from math_commands.tex:202
\providecommand{\cX}{\mathcal{X}}
% from math_commands.tex:203
\providecommand{\cY}{\mathcal{Y}}
% from math_commands.tex:204
\providecommand{\cW}{\mathcal{W}}
% from math_commands.tex:205
\providecommand{\cZ}{\mathcal{Z}}
% from math_commands.tex:207
\providecommand{\error}{\mathcal{E}}
% from math_commands.tex:208
\providecommand{\ce}{\mathcal{C}}
% from math_commands.tex:209
\providecommand{\tce}{\Xi}
% from math_commands.tex:222
\newcommand{\eg}{e.g.}
% from math_commands.tex:223
\newcommand{\ie}{i.e.}
% from math_commands.tex:224
\newcommand{\iid}{i.i.d.}
% from math_commands.tex:225
\newcommand{\cf}{cf.}
% from math_commands.tex:228
\newcommand{\ti}{\widetilde}
% from math_commands.tex:229
\newcommand{\ov}{\overline}
% from math_commands.tex:232
\newcommand{\tgamma}{\tilde \gamma}
% from math_commands.tex:233
\newcommand{\true}{\texttt{true}}
% from math_commands.tex:234
\newcommand{\false}{\texttt{false}}
% from math_commands.tex:240
\newcommand{\cmark}{\ding{51}}
% from math_commands.tex:241
\newcommand{\xmark}{\ding{55}}
% from math_commands.tex:242
\newcommand{\yes}{$\checkmark$}
% from math_commands.tex:243
\newcommand{\no}{$\times$}
% from math_commands.tex:245
\newcommand{\minus}{\scalebox{0.8}{$-$}}
% from math_commands.tex:246
\newcommand{\plus}{\scalebox{0.6}{$+$}}
% from math_commands.tex:256
\newcommand*{\tikzmk}[1]{\tikz[remember picture,overlay,] \node (#1) {};}
% from math_commands.tex:258
\newcommand{\boxit}[1]{\tikz[remember picture,overlay]{\node[yshift=3pt,fill=#1,opacity=.25,fit={($(A)+(-0.2*\linewidth - 3pt,3pt)$)($(B)+(0.75*\linewidth - 5pt,-2pt)$)}] {};}}
% from math_commands.tex:260
\newcommand{\blockfed}[1]{\tikz[remember picture,overlay]{\node[yshift=3pt,fill=#1,opacity=.25,fit={($(A)+(-0.1*\linewidth - 3pt,3pt)$)($(B)+(0.88*\linewidth - 5pt,-2pt)$)}] {};}}
% from math_commands.tex:262
\newcommand{\highlight}[1]{\tikz[remember picture,overlay]{\node[yshift=3pt,fill=#1,opacity=.25,fit={($(A)+(2pt,6pt)$)($(B)+(-3pt,-7pt)$)}] {};}}
% from math_commands.tex:265
\newcommand{\speedup}[1]{{\color{gray}(\ifdim #1 pt > 0.3pt #1\else $< #1$\fi{}$\times$)}}
% from math_commands.tex:272
\newcommand{\llangle}[1][]{\savebox{\@brx}{\(\m@th{#1\langle}\)}%
  \mathopen{\copy\@brx\mkern2mu\kern-0.9\wd\@brx\usebox{\@brx}}}
% from math_commands.tex:274
\newcommand{\rrangle}[1][]{\savebox{\@brx}{\(\m@th{#1\rangle}\)}%
  \mathclose{\copy\@brx\mkern2mu\kern-0.9\wd\@brx\usebox{\@brx}}}
% from math_commands.tex:277
\newcommand{\algopt}{\textsc{Choco-SGD}\xspace}
% from math_commands.tex:278
\newcommand{\algcons}{\textsc{Choco-Gossip}\xspace}
% from math_commands.tex:279
\newcommand{\eqncons}{\textsc{Choco-G}\xspace}
% from math_commands.tex:283
\providecommand{\lin}[1]{\ensuremath{\left\langle #1 \right\rangle}}
% from math_commands.tex:284
\providecommand{\abs}[1]{\left\lvert#1\right\rvert}
% from math_commands.tex:285
\providecommand{\norm}[1]{\left\lVert#1\right\rVert}
% [intra-paper duplicate] \refLE — also defined at math_commands.tex:86
% from math_commands.tex:287
\providecommand{\refLE}[1]{\ensuremath{\stackrel{(\ref{#1})}{\leq}}}
% [intra-paper duplicate] \refEQ — also defined at math_commands.tex:87
% from math_commands.tex:288
\providecommand{\refEQ}[1]{\ensuremath{\stackrel{(\ref{#1})}{=}}}
% [intra-paper duplicate] \refGE — also defined at math_commands.tex:88
% from math_commands.tex:289
\providecommand{\refGE}[1]{\ensuremath{\stackrel{(\ref{#1})}{\geq}}}
% [intra-paper duplicate] \refID — also defined at math_commands.tex:89
% from math_commands.tex:290
\providecommand{\refID}[1]{\ensuremath{\stackrel{(\ref{#1})}{\equiv}}}
% [intra-paper duplicate] \R — also defined at math_commands.tex:97
% from math_commands.tex:292
\providecommand{\R}{\mathbb{R}}
% [intra-paper duplicate] \N — also defined at math_commands.tex:99
% from math_commands.tex:293
\providecommand{\N}{\mathbb{N}}
% [intra-paper duplicate] \Eb — also defined at math_commands.tex:17
% from math_commands.tex:296
\providecommand{\Eb}[1]{{\mathbb E}\left[#1\right] }
% from math_commands.tex:297
\providecommand{\Ec}[2]{{\mathbb E}_{#2}\left[#1\right] }
% from math_commands.tex:298
\providecommand{\EE}[2]{{\mathbb E}_{#1}\left.#2\right. }
% from math_commands.tex:299
\providecommand{\EEb}[2]{{\mathbb E}_{#1}\left[#2\right] }
% from math_commands.tex:300
\providecommand{\prob}[1]{{\rm Pr}\left[#1\right] }
% from math_commands.tex:301
\providecommand{\Prob}[2]{{\rm Pr}_{#1}\left[#2\right] }
% from math_commands.tex:302
\providecommand{\P}[1]{{\rm Pr}\left.#1\right. }
% from math_commands.tex:303
\providecommand{\Pb}[1]{{\rm Pr}\left[#1\right] }
% from math_commands.tex:304
\providecommand{\PP}[2]{{\rm Pr}_{#1}\left[#2\right] }
% from math_commands.tex:305
\providecommand{\PPb}[2]{{\rm Pr}_{#1}\left[#2\right] }
% from math_commands.tex:307
\newcommand{\Vara}[1]{\Var\left[#1\right]}
% from math_commands.tex:308
\newcommand{\Varb}[2]{\Var_{#1}\left[#2\right]}
% from math_commands.tex:309
\newcommand{\lone}[1]{\norm{#1}_1}
% from math_commands.tex:310
\newcommand{\og}{\overline{g}}
% from math_commands.tex:311
\newcommand{\oh}{\overline{H}}
% [intra-paper duplicate] \0 — also defined at math_commands.tex:110
% from math_commands.tex:312
\providecommand{\0}{\mathbf{0}}
% [intra-paper duplicate] \1 — also defined at math_commands.tex:111
% from math_commands.tex:313
\providecommand{\1}{\mathbf{1}}
% [intra-paper duplicate] \aa — also defined at math_commands.tex:113
% from math_commands.tex:314
\renewcommand{\aa}{\mathbf{a}}
% [intra-paper duplicate] \bb — also defined at math_commands.tex:114
% from math_commands.tex:315
\providecommand{\bb}{\mathbf{b}}
% [intra-paper duplicate] \cc — also defined at math_commands.tex:115
% from math_commands.tex:316
\providecommand{\cc}{\mathbf{c}}
% [intra-paper duplicate] \dd — also defined at math_commands.tex:116
% from math_commands.tex:317
\providecommand{\dd}{\mathbf{d}}
% [intra-paper duplicate] \ee — also defined at math_commands.tex:117
% from math_commands.tex:318
\providecommand{\ee}{\mathbf{e}}
% [intra-paper duplicate] \ff — also defined at math_commands.tex:118
% from math_commands.tex:319
\providecommand{\ff}{\mathbf{f}}
% [intra-paper duplicate] \hh — also defined at math_commands.tex:120
% from math_commands.tex:320
\providecommand{\hh}{\mathbf{h}}
% [intra-paper duplicate] \ii — also defined at math_commands.tex:121
% from math_commands.tex:321
\providecommand{\ii}{\mathbf{i}}
% [intra-paper duplicate] \jj — also defined at math_commands.tex:122
% from math_commands.tex:322
\providecommand{\jj}{\mathbf{j}}
% from math_commands.tex:323
\providecommand{\erf}{\operatorname{erf}}
% [intra-paper duplicate] \lll — also defined at math_commands.tex:124
% from math_commands.tex:324
\let\lll\ll
% [intra-paper duplicate] \mm — also defined at math_commands.tex:126
% from math_commands.tex:326
\providecommand{\mm}{\mathbf{m}}
% [intra-paper duplicate] \nn — also defined at math_commands.tex:127
% from math_commands.tex:327
\providecommand{\nn}{\mathbf{n}}
% [intra-paper duplicate] \oo — also defined at math_commands.tex:128
% from math_commands.tex:328
\providecommand{\oo}{\mathbf{o}}
% [intra-paper duplicate] \pp — also defined at math_commands.tex:129
% from math_commands.tex:329
\providecommand{\pp}{\mathbf{p}}
% [intra-paper duplicate] \qq — also defined at math_commands.tex:130
% from math_commands.tex:330
\providecommand{\qq}{\mathbf{q}}
% [intra-paper duplicate] \rr — also defined at math_commands.tex:131
% from math_commands.tex:331
\providecommand{\rr}{\mathbf{r}}
% [intra-paper duplicate] \ss — also defined at math_commands.tex:133
% from math_commands.tex:332
\providecommand{\ss}{\mathbb{s}}
% [intra-paper duplicate] \tt — also defined at math_commands.tex:134
% from math_commands.tex:333
\providecommand{\tt}{\mathbf{t}}
% [intra-paper duplicate] \uu — also defined at math_commands.tex:135
% from math_commands.tex:334
\providecommand{\uu}{\mathbf{u}}
% [intra-paper duplicate] \vv — also defined at math_commands.tex:136
% from math_commands.tex:335
\providecommand{\vv}{\mathbf{v}}
% [intra-paper duplicate] \ww — also defined at math_commands.tex:137
% from math_commands.tex:336
\providecommand{\ww}{\mathbf{w}}
% [intra-paper duplicate] \xx — also defined at math_commands.tex:138
% from math_commands.tex:337
\providecommand{\xx}{\mathbf{x}}
% [intra-paper duplicate] \yy — also defined at math_commands.tex:139
% from math_commands.tex:338
\providecommand{\yy}{\mathbf{y}}
% [intra-paper duplicate] \zz — also defined at math_commands.tex:142
% from math_commands.tex:339
\providecommand{\zz}{\mathbf{z}}
% [intra-paper duplicate] \mA — also defined at math_commands.tex:151
% from math_commands.tex:341
\providecommand{\mA}{\mathbf{A}}
% [intra-paper duplicate] \mB — also defined at math_commands.tex:152
% from math_commands.tex:342
\providecommand{\mB}{\mathbf{B}}
% [intra-paper duplicate] \mC — also defined at math_commands.tex:153
% from math_commands.tex:343
\providecommand{\mC}{\mathbf{C}}
% [intra-paper duplicate] \mD — also defined at math_commands.tex:154
% from math_commands.tex:344
\providecommand{\mD}{\mathbf{D}}
% [intra-paper duplicate] \mE — also defined at math_commands.tex:155
% from math_commands.tex:345
\providecommand{\mE}{\mathbf{E}}
% [intra-paper duplicate] \mF — also defined at math_commands.tex:156
% from math_commands.tex:346
\providecommand{\mF}{\mathbf{F}}
% [intra-paper duplicate] \mG — also defined at math_commands.tex:157
% from math_commands.tex:347
\providecommand{\mG}{\mathbf{G}}
% [intra-paper duplicate] \mH — also defined at math_commands.tex:158
% from math_commands.tex:348
\providecommand{\mH}{\mathbf{H}}
% [intra-paper duplicate] \mI — also defined at math_commands.tex:159
% from math_commands.tex:349
\providecommand{\mI}{\mathbf{I}}
% [intra-paper duplicate] \mJ — also defined at math_commands.tex:160
% from math_commands.tex:350
\providecommand{\mJ}{\mathbf{J}}
% [intra-paper duplicate] \mK — also defined at math_commands.tex:161
% from math_commands.tex:351
\providecommand{\mK}{\mathbf{K}}
% [intra-paper duplicate] \mL — also defined at math_commands.tex:162
% from math_commands.tex:352
\providecommand{\mL}{\mathbf{L}}
% [intra-paper duplicate] \mM — also defined at math_commands.tex:163
% from math_commands.tex:353
\providecommand{\mM}{\mathbf{M}}
% [intra-paper duplicate] \mN — also defined at math_commands.tex:164
% from math_commands.tex:354
\providecommand{\mN}{\mathbf{N}}
% [intra-paper duplicate] \mO — also defined at math_commands.tex:165
% from math_commands.tex:355
\providecommand{\mO}{\mathbf{O}}
% [intra-paper duplicate] \mP — also defined at math_commands.tex:166
% from math_commands.tex:356
\providecommand{\mP}{\mathbf{P}}
% [intra-paper duplicate] \mQ — also defined at math_commands.tex:167
% from math_commands.tex:357
\providecommand{\mQ}{\mathbf{Q}}
% [intra-paper duplicate] \mR — also defined at math_commands.tex:168
% from math_commands.tex:358
\providecommand{\mR}{\mathbf{R}}
% [intra-paper duplicate] \mS — also defined at math_commands.tex:169
% from math_commands.tex:359
\providecommand{\mS}{\mathbf{S}}
% [intra-paper duplicate] \mT — also defined at math_commands.tex:170
% from math_commands.tex:360
\providecommand{\mT}{\mathbf{T}}
% [intra-paper duplicate] \mU — also defined at math_commands.tex:171
% from math_commands.tex:361
\providecommand{\mU}{\mathbf{U}}
% [intra-paper duplicate] \mV — also defined at math_commands.tex:172
% from math_commands.tex:362
\providecommand{\mV}{\mathbf{V}}
% [intra-paper duplicate] \mW — also defined at math_commands.tex:173
% from math_commands.tex:363
\providecommand{\mW}{\mathbf{W}}
% [intra-paper duplicate] \mX — also defined at math_commands.tex:174
% from math_commands.tex:364
\providecommand{\mX}{\mathbf{X}}
% [intra-paper duplicate] \mY — also defined at math_commands.tex:175
% from math_commands.tex:365
\providecommand{\mY}{\mathbf{Y}}
% [intra-paper duplicate] \mZ — also defined at math_commands.tex:176
% from math_commands.tex:366
\providecommand{\mZ}{\mathbf{Z}}
% [intra-paper duplicate] \mLambda — also defined at math_commands.tex:177
% from math_commands.tex:367
\providecommand{\mLambda}{\mathbf{\Lambda}}
% from math_commands.tex:368
\providecommand{\mlambda}{\bm{\lambda}}
% from math_commands.tex:369
\providecommand{\mmu}{\bm{\mu}}
% [intra-paper duplicate] \cA — also defined at math_commands.tex:180
% from math_commands.tex:370
\providecommand{\cA}{\mathcal{A}}
% [intra-paper duplicate] \cB — also defined at math_commands.tex:181
% from math_commands.tex:371
\providecommand{\cB}{\mathcal{B}}
% [intra-paper duplicate] \cC — also defined at math_commands.tex:182
% from math_commands.tex:372
\providecommand{\cC}{\mathcal{C}}
% [intra-paper duplicate] \cD — also defined at math_commands.tex:183
% from math_commands.tex:373
\providecommand{\cD}{\mathcal{D}}
% [intra-paper duplicate] \cE — also defined at math_commands.tex:184
% from math_commands.tex:374
\providecommand{\cE}{\mathcal{E}}
% [intra-paper duplicate] \cF — also defined at math_commands.tex:185
% from math_commands.tex:375
\providecommand{\cF}{\mathcal{F}}
% [intra-paper duplicate] \cG — also defined at math_commands.tex:186
% from math_commands.tex:376
\providecommand{\cG}{\mathcal{G}}
% [intra-paper duplicate] \cH — also defined at math_commands.tex:187
% from math_commands.tex:377
\providecommand{\cH}{\mathcal{H}}
% [intra-paper duplicate] \cJ — also defined at math_commands.tex:189
% from math_commands.tex:378
\providecommand{\cJ}{\mathcal{J}}
% [intra-paper duplicate] \cK — also defined at math_commands.tex:190
% from math_commands.tex:379
\providecommand{\cK}{\mathcal{K}}
% [intra-paper duplicate] \cL — also defined at math_commands.tex:191
% from math_commands.tex:380
\providecommand{\cL}{\mathcal{L}}
% [intra-paper duplicate] \cM — also defined at math_commands.tex:192
% from math_commands.tex:381
\providecommand{\cM}{\mathcal{M}}
% [intra-paper duplicate] \cN — also defined at math_commands.tex:193
% from math_commands.tex:382
\providecommand{\cN}{\mathcal{N}}
% [intra-paper duplicate] \cO — also defined at math_commands.tex:194
% from math_commands.tex:383
\providecommand{\cO}{\mathcal{O}}
% [intra-paper duplicate] \cP — also defined at math_commands.tex:195
% from math_commands.tex:384
\providecommand{\cP}{\mathcal{P}}
% [intra-paper duplicate] \cQ — also defined at math_commands.tex:196
% from math_commands.tex:385
\providecommand{\cQ}{\mathcal{Q}}
% [intra-paper duplicate] \cR — also defined at math_commands.tex:197
% from math_commands.tex:386
\providecommand{\cR}{\mathcal{R}}
% [intra-paper duplicate] \cS — also defined at math_commands.tex:80, math_commands.tex:198
% from math_commands.tex:387
\providecommand{\cS}{\mathcal{S}}
% [intra-paper duplicate] \cT — also defined at math_commands.tex:199
% from math_commands.tex:388
\providecommand{\cT}{\mathcal{T}}
% [intra-paper duplicate] \cU — also defined at math_commands.tex:200
% from math_commands.tex:389
\providecommand{\cU}{\mathcal{U}}
% [intra-paper duplicate] \cV — also defined at math_commands.tex:201
% from math_commands.tex:390
\providecommand{\cV}{\mathcal{V}}
% [intra-paper duplicate] \cX — also defined at math_commands.tex:202
% from math_commands.tex:391
\providecommand{\cX}{\mathcal{X}}
% [intra-paper duplicate] \cY — also defined at math_commands.tex:203
% from math_commands.tex:392
\providecommand{\cY}{\mathcal{Y}}
% [intra-paper duplicate] \cW — also defined at math_commands.tex:204
% from math_commands.tex:393
\providecommand{\cW}{\mathcal{W}}
% [intra-paper duplicate] \cZ — also defined at math_commands.tex:205
% from math_commands.tex:394
\providecommand{\cZ}{\mathcal{Z}}
% [intra-paper duplicate] \vec — also defined at assembler.tex:29, math_commands.tex:85
% from math_commands.tex:396
\renewcommand{\vec}{\mathbf}
% from math_commands.tex:399
\providecommand{\mycomment}[3]{\todo[caption={},size=footnotesize,color=#1!20]{\textbf{#2: }#3}}
% from math_commands.tex:400
\providecommand{\inlinecomment}[3]{%
  {\color{#1}#2: #3}}
% from math_commands.tex:402
\newcommand
% from assembler.tex:67
\newcommand{\vDelta}{\mathbf{\Delta}}
% from assembler.tex:68
\def\hR{\widehat{\mathcal{R}}}
% from assembler.tex:69
\newcommand

% --- Environments ---
% from assembler.tex:71
\newenvironment{assumptionp}[1]{
  \renewcommand\theassumptionalt{#1}
  \assumptionalt
}{\endassumptionalt}

% --- Theorems ---
% from assembler.tex:47
\newtheorem*{algorithm*}{Algorithm}
% from assembler.tex:48
\newtheorem{algorithm}{Algorithm}
% from assembler.tex:49
\newtheorem{assumption}{Assumption}
% from assembler.tex:50
\newtheorem{myalgorithm*}{Algorithm}
% from assembler.tex:51
\newtheorem{myalgorithm}{Algorithm}
% from assembler.tex:52
\newtheorem{theorem}{Theorem}
% from assembler.tex:53
\newtheorem{lemma}{Lemma}
% from assembler.tex:54
\newtheorem{proposition}{Proposition}
% from assembler.tex:55
\newtheorem{definition}{Definition}
% from assembler.tex:56
\newtheorem{corollary}{Corollary}
% from assembler.tex:57
\newtheorem{conjecture}{Conjecture}
% from assembler.tex:58
\newtheorem{remark}{Remark}
% from assembler.tex:59
\newtheorem{example}{Example}
